%=====================================================================
\chapter{Semi-Supervised Learning and the EM Algorithm}\label{ch:semisup}
%=====================================================================
\locator{4}{Part IV --- Learning Without Full Labels}

\begin{outcomes}
By the end of this chapter you will be able to:
\begin{tight}
  \item \textbf{Explain} when unlabelled examples can improve a classifier, and the
        assumptions that make it possible.
  \item \textbf{Describe} self-training, and \textbf{state} its main risk.
  \item \textbf{Formulate} a Gaussian mixture model with hidden cluster labels.
  \item \textbf{Perform} an iteration of the EM algorithm by hand: responsibilities
        in the E-step, weighted means in the M-step.
  \item \textbf{Apply} EM to a mixture of labelled and unlabelled data, and
        \textbf{relate} $k$-means to EM.
\end{tight}
\end{outcomes}

\begin{prereq}
\chref{math} (the Gaussian, maximum likelihood), \chref{bayes} (posteriors and
naive Bayes) and \chref{clustering} ($k$-means). The algorithm of this chapter
returns in \chref{hmm}.
\end{prereq}

\begin{keyterms}
semi-supervised learning $\bullet$ smoothness, cluster and manifold assumptions
$\bullet$ self-training (pseudo-labelling) $\bullet$ confirmation bias $\bullet$
latent (hidden) variable $\bullet$ mixture model $\bullet$ Gaussian mixture
$\bullet$ mixing weight $\bullet$ responsibility $\bullet$ EM algorithm
$\bullet$ E-step $\bullet$ M-step $\bullet$ local maximum $\bullet$ hard vs.\ soft
assignment $\bullet$ label propagation
\end{keyterms}

\section{Few Labels, Many Examples}

Labels are expensive. A radiologist labels a scan; a security analyst investigates a
binary; a teacher marks an essay. Unlabelled examples are nearly free: every scan in
the archive, every file on every machine, every essay ever submitted.
\term{Semi-supervised learning} uses both --- a small labelled set and a large
unlabelled one --- to learn a better classifier than the labelled set alone would
give.

How can examples with no labels possibly help? Only if the way the examples are
spread out says something about where the classes lie. Three assumptions, all
versions of the same idea, make that true:
\begin{tight}
  \item \textbf{Smoothness:} examples close together are likely to share a label.
  \item \textbf{Cluster:} examples form clusters, and a cluster tends to be one
        class --- so the decision boundary should pass through low-density gaps,
        not through the middle of a cluster.
  \item \textbf{Manifold:} high-dimensional data lies near a lower-dimensional
        surface, and distances along that surface are what matter.
\end{tight}

\dsfig{%
\begin{tikzpicture}[font=\scriptsize,
  unl/.style={circle, fill=gray!40, inner sep=1.2pt},
  la/.style={circle, fill=secondary, inner sep=2.6pt, draw=white, line width=0.6pt},
  lb/.style={rectangle, fill=accent, inner sep=2.8pt, draw=white, line width=0.6pt}]
\foreach \ox/\ttl in {0/{Two labels only}, 7.2/{Two labels + unlabelled data}}{
  \begin{scope}[shift={(\ox,0)}]
  \draw[rounded corners=4pt, gray!40] (-0.2,-0.2) rectangle (6.2,4.2);
  \node[font=\scriptsize\bfseries, text=primary] at (3.0,4.6) {\ttl};
  \foreach \x/\y in {0.8/1.9,1.1/2.6,1.4/1.4,1.7/2.2,1.2/3.1,2.0/2.9,0.6/2.4,1.9/1.6,2.3/2.4,1.5/3.4,
                     0.9/1.2,2.4/3.2,2.1/1.1,1.0/3.6,2.6/1.8}{\node[unl] at (\x,\y) {};}
  \foreach \x/\y in {4.1/2.2,4.4/1.4,4.7/2.8,5.0/1.9,4.2/3.2,5.3/2.5,4.8/1.1,5.5/3.2,3.9/1.6,
                     5.1/3.6,4.5/0.8,5.6/1.6,4.3/2.6,5.4/0.9,3.8/3.0}{\node[unl] at (\x,\y) {};}
  \node[la] at (2.3,2.4) {}; \node[lb] at (3.9,1.6) {};
  \end{scope}}
\draw[primary, line width=1.2pt, dashed] (3.62,0) -- (2.58,4.0);
\node[font=\tiny, primary, anchor=north, align=center] at (3.1,-0.25)
     {halfway between the two labels:\\cuts through the left cluster};
\draw[greenhl!70!black, line width=1.4pt] (7.2+3.2,0) -- (7.2+3.2,4.0);
\node[font=\tiny, greenhl!45!black, anchor=north, align=center] at (7.2+3.2,-0.25)
     {through the empty gap:\\each cluster one class};
\end{tikzpicture}}%
{The cluster assumption. With two labelled points (blue circle, red square) a classifier
puts the boundary between them, through the middle of a cluster. The unlabelled points
(grey) reveal two clusters and a gap, and the boundary belongs in the gap.}{cluster-assumption}

\begin{alertbox}[Unlabelled Data Can Also Hurt]
When the assumptions are false --- the classes overlap in one blob, or a cluster
contains both classes --- unlabelled data pushes the boundary to the wrong place,
and the semi-supervised model can be \emph{worse} than one trained on the labels
alone. Always compare against the labelled-only baseline on a held-out labelled test
set.
\end{alertbox}

\section{Self-Training}

The simplest method needs no new theory. \term{Self-training} (pseudo-labelling):
\begin{tightnum}
  \item train a classifier on the labelled examples;
  \item predict the unlabelled examples, and add those predicted with high
        confidence to the training set, with their predicted labels;
  \item retrain, and repeat until no confident predictions remain.
\end{tightnum}
It works with any classifier that gives probabilities, and on the handwritten digits
of this chapter's lab it lifts accuracy from 80\% with 50 labels to 90\% using the
1{,}207 unlabelled images as well (all 1{,}257 labels give 97\%). Its risk is
\term{confirmation bias}: an early confident mistake becomes a training label, is
reinforced at every round, and is never corrected. A high confidence threshold
limits the damage.

\section{Mixture Models and Hidden Variables}

A more principled approach models how the data was generated. Suppose each example
was produced by first picking a class (or cluster) $z$, then drawing $\mathbf{x}$
from that class's distribution. For labelled examples $z$ is observed; for
unlabelled ones it is a \term{hidden} (latent) variable.

\begin{definitionbox}[Gaussian Mixture Model]
A \term{mixture} of $k$ Gaussians has, for each component $j$, a \term{mixing
weight} $\pi_j$ (the probability of choosing it, with $\sum_j \pi_j = 1$), a mean
$\boldsymbol{\mu}_j$ and a covariance $\Sigma_j$. The probability density of an example
is
\[
p(\mathbf{x}) = \sum_{j=1}^{k} \pi_j\,\mathcal{N}(\mathbf{x} \mid \boldsymbol{\mu}_j, \Sigma_j)
\]
If we knew which component produced each example, fitting would be easy: average
each component's own examples. If we knew the parameters, we could compute for each
example the probability that each component produced it. We know neither --- and
the EM algorithm alternates between the two.
\end{definitionbox}

\section{The EM Algorithm}

\begin{definitionbox}[Expectation--Maximisation (EM)]
Start from initial parameter guesses. Repeat until the likelihood stops improving:
\begin{tight}
  \item \textbf{E-step:} with the current parameters, compute for every example
        $i$ and component $j$ the \term{responsibility} $r_{ij}$ below --- the
        posterior probability that component $j$ produced example $i$.
  \item \textbf{M-step:} re-estimate each component from \emph{all} examples,
        weighted by responsibility: $\boldsymbol{\mu}_j = \sum_i r_{ij}\mathbf{x}_i / \sum_i r_{ij}$,
        $\pi_j = \tfrac{1}{n}\sum_i r_{ij}$, and $\Sigma_j$ likewise.
\end{tight}
\[
r_{ij} = \frac{\pi_j\,\mathcal{N}(\mathbf{x}_i \mid \boldsymbol{\mu}_j, \Sigma_j)}
              {\sum_{l}\pi_l\,\mathcal{N}(\mathbf{x}_i \mid \boldsymbol{\mu}_l, \Sigma_l)}
\]
Each iteration never decreases the likelihood of the data, so EM converges --- to a
\term{local maximum}, which depends on the starting point.
\end{definitionbox}

\begin{worked}[EM for Two Means (Mitchell's Example, in Numbers)]
Six values: $x = 1, 2, 3, 7, 8, 9$. Model: two Gaussians with known $\sigma = 1$ and equal
weights; unknown means. Start from $\mu_1 = 2$, $\mu_2 = 5$.

\smallskip
\textbf{E-step.} For each $x$, the densities are proportional to
$e^{-(x-\mu)^2/2}$:
\begin{tight}
  \item $x = 1$: $e^{-0.5} = 0.607$ against $e^{-8} = 0.0003$ $\Rightarrow$ $r_1 = 0.999$
  \item $x = 2$: $e^{0} = 1$ against $e^{-4.5} = 0.011$ $\Rightarrow$ $r_1 = 0.989$
  \item $x = 3$: $e^{-0.5} = 0.607$ against $e^{-2} = 0.135$ $\Rightarrow$ $r_1 = 0.818$
  \item $x = 7, 8, 9$: component 2 is overwhelmingly likelier $\Rightarrow$ $r_1 \approx 0$
\end{tight}
\textbf{M-step.} Responsibility-weighted means:
\[
\mu_1 = \frac{0.999(1) + 0.989(2) + 0.818(3)}{0.999 + 0.989 + 0.818} = \frac{5.430}{2.806} = \mathbf{1.935}
\]
\[
\mu_2 = \frac{0.001(1) + 0.011(2) + 0.182(3) + 7 + 8 + 9}{0.001 + 0.011 + 0.182 + 3} = \frac{24.570}{3.194} = \mathbf{7.693}
\]
Note that $x = 3$ contributes 18\% of itself to component 2: a \emph{soft}
assignment.

\smallskip
\textbf{Iteration 2} moves the means to $2.000$ and $8.000$, after which nothing
changes. The log-likelihood of the data rises from $-24.96$ at the starting means to
$-11.82$ after one iteration and $-11.67$ at convergence. EM has found the two groups
\{1, 2, 3\} and \{7, 8, 9\}.
\end{worked}

\dsfig{%
\begin{tikzpicture}
\begin{groupplot}[
  group style={group size=2 by 1, horizontal sep=1.3cm},
  width=6.9cm, height=4.4cm, axis lines=left,
  tick label style={font=\tiny}, label style={font=\tiny},
  title style={font=\scriptsize\bfseries}, xlabel={$x$},
  xmin=-1, xmax=11, ymin=0, ymax=0.46, ytick=\empty]
\nextgroupplot[title={Start: $\mu_1 = 2$, $\mu_2 = 5$}]
\addplot[secondary, line width=1.2pt, domain=-1:11, samples=120] {0.5*exp(-(x-2)^2/2)/sqrt(2*pi)*2};
\addplot[accent, line width=1.2pt, domain=-1:11, samples=120] {0.5*exp(-(x-5)^2/2)/sqrt(2*pi)*2};
\addplot[only marks, mark=*, mark size=2.2pt, secondary] coordinates {(1,0.02)(2,0.02)};
\addplot[only marks, mark=*, mark size=2.2pt, purplehl] coordinates {(3,0.02)};
\addplot[only marks, mark=*, mark size=2.2pt, accent] coordinates {(7,0.02)(8,0.02)(9,0.02)};
\draw[purplehl, line width=0.4pt] (axis cs:3,0.04) -- (axis cs:3.4,0.36);
\node[font=\tiny, purplehl, anchor=south west] at (axis cs:3.3,0.36) {$x = 3$: 82\% / 18\%};
\nextgroupplot[title={After EM: $\mu_1 = 2$, $\mu_2 = 8$}]
\addplot[secondary, line width=1.2pt, domain=-1:11, samples=120] {0.5*exp(-(x-2)^2/2)/sqrt(2*pi)*2};
\addplot[accent, line width=1.2pt, domain=-1:11, samples=120] {0.5*exp(-(x-8)^2/2)/sqrt(2*pi)*2};
\addplot[only marks, mark=*, mark size=2.2pt, secondary] coordinates {(1,0.02)(2,0.02)(3,0.02)};
\addplot[only marks, mark=*, mark size=2.2pt, accent] coordinates {(7,0.02)(8,0.02)(9,0.02)};
\end{groupplot}
\end{tikzpicture}}%
{EM on the worked example. Each point is coloured by its responsibilities; $x = 3$ starts
out shared between the components. Two iterations move the second Gaussian onto the group
it explains.}{em-1d}

\begin{conceptbox}[$k$-Means Is EM With Hard Assignments]
Replace the E-step's probabilities by a hard choice --- responsibility 1 for the
nearest mean, 0 for the others --- and the M-step becomes ``each mean is the average
of its assigned points''. That is exactly $k$-means (\chref{clustering}). EM for
Gaussian mixtures is the soft, probabilistic generalisation: clusters may overlap,
may have different sizes and elongated shapes (through $\Sigma_j$), and each example
gets a probability of membership instead of a verdict.
\end{conceptbox}

\dsfig{%
\begin{tikzpicture}[font=\scriptsize,
  st/.style={rounded corners=4pt, draw=#1, line width=1pt, fill=#1!9, text=#1!55!black,
             align=center, minimum width=4.4cm, minimum height=1.4cm, font=\scriptsize},
  a/.style={-{Stealth[length=2.2mm]}, gray!60, line width=1pt}]
\node[st=secondary] (e) at (0,0) {\textbf{E-step}\\[2pt]fill in the hidden labels\\as probabilities $r_{ij}$};
\node[st=greenhl!70!black] (m) at (6.4,0) {\textbf{M-step}\\[2pt]re-fit the parameters\\as if the $r_{ij}$ were counts};
\draw[a] (e.north east) .. controls +(1.4,0.9) and +(-1.4,0.9) .. node[above, font=\tiny] {soft labels} (m.north west);
\draw[a] (m.south west) .. controls +(-1.4,-0.9) and +(1.4,-0.9) .. node[below, font=\tiny] {new parameters} (e.south east);
\node[anchor=west, align=left, font=\scriptsize, text=gray!55!black] at (9.0,0)
     {The same two steps fit:\\[2pt]
      \textbullet\ Gaussian mixtures\\
      \textbullet\ naive Bayes plus unlabelled data\\
      \textbullet\ hidden Markov models (Ch.\ 18)\\
      \textbullet\ $k$-means (hard assignments)};
\end{tikzpicture}}%
{The EM loop. Whenever a model would be easy to fit if some variable were observed, EM
fills that variable in with probabilities and fits as if it were observed.}{em-loop}

\section{EM With Labelled and Unlabelled Data}

This is the method the HEC outline names. Nigam, McCallum, Thrun and Mitchell (2000)
used it to classify web pages and news articles from a handful of labelled documents
and thousands of unlabelled ones. The recipe:
\begin{tightnum}
  \item Fit the model --- naive Bayes for text, or a Gaussian per class --- on the
        labelled examples alone.
  \item \textbf{E-step:} give each \emph{unlabelled} example a probability for each
        class from the current model. Labelled examples keep their known labels
        (responsibility 1 for their class, 0 for the others).
  \item \textbf{M-step:} refit the model on all examples, labelled ones counting in
        full and unlabelled ones fractionally, by their responsibilities.
  \item Repeat until the parameters settle.
\end{tightnum}
The labelled examples anchor each component to a class --- they say \emph{which}
cluster is which --- and the unlabelled examples say \emph{where} the clusters really
are. The authors report that adding unlabelled documents to small labelled sets
reduced text-classification error by up to 30\%.

\begin{worked}[Two Labels and Four Unlabelled Points]
The six values of the previous example, but now $x = 3$ is labelled class A and $x = 7$
class B; the rest are unlabelled.

\smallskip
\textbf{Labelled data alone} gives $\mu_A = 3$, $\mu_B = 7$, and a boundary at 5.

\textbf{E-step.} $x = 1$ and $x = 2$ are far likelier under A, and $x = 8$ and $x = 9$
under B; $x = 3$ and $x = 7$ are clamped to their labels.

\textbf{M-step.} $\mu_A = (1 + 2 + 3)/3 = 2$ and $\mu_B = (7 + 8 + 9)/3 = 8$, to
within a thousandth. The next iteration changes nothing.

\smallskip
The class means moved from 3 and 7 to \textbf{2 and 8} using points nobody labelled.
Here the boundary stays at 5, but it would not if the two groups had different
spreads or sizes --- and in higher dimensions the direction of the boundary moves too.
\end{worked}

\textbf{Graph-based methods} such as \term{label propagation} take the smoothness
assumption literally: connect each example to its nearest neighbours in a graph, fix
the labelled nodes, and let labels diffuse along the edges until they settle. They
follow curved clusters that a Gaussian mixture cannot represent.

\begin{pitfall}
\begin{tight}
  \item \textbf{Skipping the labelled-only baseline.} Unlabelled data can make a
        model worse; prove that it helped.
  \item \textbf{Evaluating on pseudo-labels.} Test only on examples with genuine
        labels, set aside before training.
  \item \textbf{A low self-training threshold.} Early mistakes become permanent
        labels.
  \item \textbf{Running EM once.} It finds a local maximum; start it several times
        and keep the highest likelihood.
  \item \textbf{A mixture with the wrong shape.} If the classes are not roughly
        Gaussian clusters, the mixture's components will not line up with them.
\end{tight}
\end{pitfall}

%---------------------------------------------------------------------
\begin{fourlenses}{Semi-Supervised Learning}
\lensLA{A few hundred essays have been marked by a teacher; tens of thousands have
not. Self-training a quality classifier on both is attractive, and dangerous: an
early confident mistake about a style of writing becomes a rule. Keep a human in the
loop for the pseudo-labels that change decisions.}

\lensPM{Only a few projects have a formal post-mortem label (``failed because of
scope creep''), but all projects have metrics. EM over labelled and unlabelled
projects estimates the typical metric profile of each failure mode from far more than
the labelled handful.}

\lensIOT{Engineers can label only a few fault episodes, but sensors record years of
unlabelled operation. A Gaussian mixture of normal operating modes, anchored by the
few labelled faults, is a practical starting point for a fault detector.}

\lensSEC{Security labels are scarce, slow and expensive --- each needs an analyst.
Semi-supervised malware and intrusion classifiers are standard, and so is their
failure mode: an attacker who knows the model self-trains can feed it examples
designed to be pseudo-labelled benign.}
\end{fourlenses}

%---------------------------------------------------------------------
\begin{lab}[EM by Hand, Then Fifty Labels and a Thousand Guesses]
\begin{lstlisting}[language=Python]
import numpy as np
from sklearn.datasets import load_digits
from sklearn.model_selection import train_test_split
from sklearn.linear_model import LogisticRegression
from sklearn.semi_supervised import SelfTrainingClassifier
from sklearn.pipeline import make_pipeline
from sklearn.preprocessing import StandardScaler

# --- 1. EM for two Gaussian means (sigma = 1 known, equal weights) -------
x = np.array([1.0, 2.0, 3.0, 7.0, 8.0, 9.0])
mu = np.array([2.0, 5.0])                        # initial guesses
for it in range(1, 5):
    dens = np.exp(-0.5 * (x[:, None] - mu[None, :]) ** 2)   # E-step
    r = dens / dens.sum(axis=1, keepdims=True)    # responsibilities
    ll = np.log(dens.mean(axis=1) / np.sqrt(2 * np.pi)).sum()
    mu = (r * x[:, None]).sum(axis=0) / r.sum(axis=0)       # M-step
    print(f"iteration {it}: log-likelihood {ll:8.3f} -> mu = {mu.round(3)}")

# --- 2. two labels, four unlabelled points -------------------------------
labelled = {2: 0, 3: 1}                 # x=3 is class A (0), x=7 class B (1)
mu = np.array([x[2], x[3]])             # labelled points alone: 3 and 7
print("from the two labels alone: mu =", mu)
for it in range(10):
    dens = np.exp(-0.5 * (x[:, None] - mu[None, :]) ** 2)
    r = dens / dens.sum(axis=1, keepdims=True)
    for i, c in labelled.items():       # a labelled point keeps its label
        r[i] = np.eye(2)[c]
    mu = (r * x[:, None]).sum(axis=0) / r.sum(axis=0)
print("EM with the unlabelled points too: mu =", mu.round(3))

# --- 3. self-training with 50 labels out of 1,257 ------------------------
X, y = load_digits(return_X_y=True)
X_tr, X_te, y_tr, y_te = train_test_split(X, y, test_size=0.3,
                                          random_state=0, stratify=y)
rng = np.random.default_rng(0)
keep = rng.choice(len(y_tr), size=50, replace=False)
y_part = np.full_like(y_tr, -1)                 # -1 means "unlabelled"
y_part[keep] = y_tr[keep]
base = make_pipeline(StandardScaler(), LogisticRegression(max_iter=2000))
only = base.fit(X_tr[keep], y_tr[keep]).score(X_te, y_te)
self_t = SelfTrainingClassifier(
    make_pipeline(StandardScaler(), LogisticRegression(max_iter=2000)),
    threshold=0.9).fit(X_tr, y_part)
print(f"50 labels only: {only:.3f}   self-training with the rest "
      f"unlabelled: {self_t.score(X_te, y_te):.3f}   "
      f"all {len(y_tr)} labels: "
      f"{base.fit(X_tr, y_tr).score(X_te, y_te):.3f}")
\end{lstlisting}
Change the self-training \texttt{threshold} to 0.5 and then 0.99. Which is better, and
what does the answer say about confirmation bias?
\end{lab}

%---------------------------------------------------------------------
\begin{checkpoint}
\begin{tightnum}
  \item State the cluster assumption, and give an example where it fails.
  \item In one EM iteration, what does the E-step compute and what does the M-step
        compute?
  \item How does $k$-means differ from EM for a Gaussian mixture?
\end{tightnum}
\end{checkpoint}

%---------------------------------------------------------------------
\begin{chaptersummary}
\begin{tight}
  \item Semi-supervised learning uses many unlabelled examples to improve a
        classifier trained on few labels, relying on the smoothness, cluster or
        manifold assumption; when those fail it can hurt.
  \item Self-training adds confident predictions as labels and retrains; its risk is
        confirmation bias.
  \item A Gaussian mixture explains data as coming from hidden components, each with
        a weight, mean and covariance.
  \item EM alternates an E-step (responsibilities: soft assignments) and an M-step
        (responsibility-weighted re-estimates), never lowering the likelihood and
        converging to a local maximum; $k$-means is its hard-assignment special case.
  \item With labelled and unlabelled data, labelled examples keep their labels in
        the E-step and anchor each component to a class.
\end{tight}
\end{chaptersummary}

\begin{reviewq}
  \item \textbf{(MCQ)} In the EM algorithm, the responsibilities are computed in:
        (a)~the M-step (b)~the E-step (c)~the initialisation (d)~neither step.
  \item \textbf{(MCQ)} $k$-means corresponds to EM with: (a)~soft assignments
        (b)~hard assignments (c)~no M-step (d)~labelled data only.
  \item \textbf{(MCQ)} Self-training's main risk is: (a)~underfitting (b)~slow
        prediction (c)~confirmation bias (d)~the curse of dimensionality.
  \item \textbf{(Short)} Why can unlabelled data improve a classifier at all?
  \item \textbf{(Short)} Explain why EM can converge to different answers from
        different starting points, and what to do about it.
  \item \textbf{(Short)} In semi-supervised EM, what role do the labelled examples
        play that the unlabelled ones cannot?
  \item \textbf{(Applied)} Repeat the first EM iteration of the worked example with
        initial means $\mu_1 = 4$, $\mu_2 = 6$. Where do the means go?
  \item \textbf{(Applied)} For the values $0, 1, 5, 6$ with $\sigma = 1$ and equal
        weights, start EM from $\mu_1 = 0$, $\mu_2 = 1$ and perform two iterations.
        Comment on the result.
\end{reviewq}
